\documentclass[10pt,a4paper]{amsart}
\usepackage[margin=19mm]{geometry}
\usepackage{amsmath,amssymb,mathtools}
\usepackage[scr=rsfs,cal=euler]{mathalfa}
\usepackage{xfrac}
\usepackage[unicode,hidelinks,bookmarks=false]{hyperref}
\newcommand{\ie}{i.e.\ }
\newcommand{\cf}{cf.\ }
\newcommand{\resp}{respectively\ }
\newcommand{\Subsection}{\subsection}
\providecommand{\nobreakdash}{\nobreak\hbox{-}}
\setlength{\emergencystretch}{2em}
\multlinegap0pt
\usepackage{bm}
\usepackage{bbm}
\usepackage{dsfont}
\usepackage{xspace}
\usepackage{cancel}
\usepackage{mathtools}
\usepackage{amsthm}
\makeatletter
\@ifundefined{thm}{\newtheorem{thm}{Thm}}{}
\@ifundefined{prop}{\newtheorem{prop}[thm]{Proposition}}{}
\makeatother
\newcommand{\ZZ}{\mathbb Z}
\newcommand{\sett}[2]{ \left\{ #1 \, \, || \, \, #2 \right \} }
\title[Skirting the $n$-tuples]{Skirting the $n$-tuples}
\author{Sam Adriaensen, Ferdinand Ihringer, William J. Martin, Ralihe R. Villagr\'an}
\date{Source: arXiv:2602.01080v1 (CC BY 4.0)}
\begin{document}
\maketitle

\noindent\emph{Source and licence.} Skirting the $n$-tuples. Version arXiv:2602.01080v1 (CC BY 4.0), distributed under the Creative Commons Attribution 4.0 International licence, as recorded in the arXiv licence metadata for this exact version and shown on the abstract page https://arxiv.org/abs/2602.01080v1. Full rights notice: licenses/arxiv-2602.01080-CC-BY-4.0.txt.\par\smallskip

\noindent\textit{Editorial scope.} Counted unit: the Prop whose source label is \texttt{None} in the article (source0.tex lines 197--202, source label \texttt{None}), delivered together with the article's own complete original proof of it (source0.tex lines 204--235, one \texttt{proof} environment of 32 lines). No mathematics is written, repaired or re-derived: statement and proof are the article's own text line for line. Required context retained uncounted: none. Every definition, display and label of those lines is retained unchanged. Omitted from this package and not redistributed: 1--196, 203--203, 236--444. Nothing inside a retained excerpt is omitted or altered: every retained body is delivered exactly as the article's lines are written, so no omission receipt of any kind accompanies this package. Citations: the retained excerpts carry no citation command at all, so no bibliography entry of the article is transcribed and no reference is redistributed. Reference adaptation: none was needed and none was made; every reference inside the delivered unit resolves inside it, so no unresolved reference or citation is produced. Printed slips kept literal: no printed word, symbol, hypothesis or number of the counted statement or of its proof was altered. Layout and class: the article's own class is \texttt{article} with packages including enumerate, fullpage, hyperref, bm, amsmath, amssymb, amsthm, url, graphicx, rotating, tikz, tkz-graph, relsize, mathrsfs; the wrapper is the lane's pinned \texttt{amsart} editorial wrapper at 10pt on A4, which restates the theorem-like environments the retained text needs and adds the article's own macro definitions that the retained text needs (\texttt{\textbackslash ZZ}, \texttt{\textbackslash sett}), with extra wrapper packages (bm, bbm, dsfont, xspace, cancel, mathtools). The delivered numbering is the unit's own. Assets: no retained span contains a figure environment, a graphics-inclusion command, a PDF-page inclusion, a table or a TikZ picture; no image file is copied into this package, the package declares no asset dependency and no image file is redistributed with it. Licence: version arXiv:2602.01080v1 of this article is distributed under the Creative Commons Attribution 4.0 International licence, as recorded in the arXiv licence metadata for this exact version and shown on the abstract page https://arxiv.org/abs/2602.01080v1. A full rights notice is delivered at licenses/arxiv-2602.01080-CC-BY-4.0.txt. Characters: every retained excerpt is pure ASCII, so the delivered engine, XeLaTeX with its default Latin Modern text font, reports no missing character.
\begin{prop}
 If $n$ is even, then
 \[
  f(n,3) \leq 2^{n/2} + \sum_{i=1}^{\lceil n/4 \rceil } \binom{n/2}{2i-1} 2^{(n/2) - 2i + 1}.
 \]
\end{prop}

\begin{proof}
Let $n$ be even.
Write $\tilde 0 = 01$, $\tilde 1 = 10$, $\bar 0 = 00$, $\bar 1 = 11$, $\bar 2 = 22$.
Define
\begin{align*}
 S_1 = \left\{ \tilde 0, \tilde 1 \right \}^{\frac n2}, &&
 S_2 = \sett{ x \in \left \{ \bar 0, \bar 1, \bar 2 \right \}^{\frac n2} }{ \text{the number of $i$'s with $x(2i) = 2$ is odd}}.
\end{align*}
We verify that $S = S_1 \cup S_2$ is a skirting set of $\ZZ_3^n$.
Take a vector $y$ in $\ZZ_3^n$, and view $y$ as a vector of $n/2$ blocks of length 2.
If none of the blocks occuring in $y$ are of the form $\bar 0$ or $\bar 1$, then $y$ is skirted by $S_1$.
Indeed, every block that occurs in $y$ either belongs to $\{10, 12, 20, 22 \}$ in which case it is skirted by $\tilde 0$, or it belongs to $\{01, 02, 21 \}$, in which case it is skirted by $\tilde 1$.
Thus, we can make an $n$-tuple $x \in S_1$ out of blocks of the form $\tilde 0$ and $\tilde 1$ that skirts $y$.

Now suppose that $y$ has a block of the form $\bar 0$ or $\bar 1$, say in the $i$th block.
Then we verify that $S_2$ skirts $y$.
Indeed, every block of length 2 contains at most 2 distinct symbols, hence it is skirted by one of the blocks $\bar 0, \bar 1, \bar 2$.
Thus, we can make a vector $x$ consisting of the blocks $\bar 0, \bar 1, \bar 2$ that skirts $y$.
If $x$ has an even number of blocks equal to $\bar 2$, then we can fix this by altering the $i$th block as follows.
Consider first the case where the $i$th block of $y$ is $\bar 0$.
Then the $i$th block of $x$ is either $\bar 1$, in which case we switch it to $\bar 2$, or it is $\bar 2$, in which case we switch it to $\bar 1$.
The resulting alteration of $x$ then has an odd number of occurrences of $\bar 2$, and hence is an element of $S_2$ that skirts $y$.
The case where the $i$th block of $y$ is $\bar 1$ is completely analogous.

We conclude that $S$ is a skirting set of $\ZZ_3^n$.
Next, we compute its size.
 The size of $S_1$ is clearly $2^{n/2}$.
 For a given number $i$, we find $\binom{n/2}{i} \cdot 2^{n/2-i}$
 elements of $\{ \bar{0}, \bar{1}, \bar{2}\}$ 
 with precisely $i$ entries equal to $\bar{2}$.
 Thus, the size of $S_1 \cup S_2$ is as asserted.
\end{proof}

\end{document}
